\documentclass[11pt]{article}
\usepackage[margin=1in]{geometry}
\usepackage{amsmath,amssymb,amsthm,booktabs,hyperref}
\hypersetup{colorlinks=true,linkcolor=blue,urlcolor=blue}
\newcommand{\sech}{\operatorname{sech}}
\newcommand{\Res}{\operatorname{Res}}
\newcommand{\Q}{\mathbb Q}
\title{The sixth-power Berndt integral\\MathDB 361364}
\author{Research note}
\date{17 September 2026}
\begin{document}
\maketitle

\section*{Task 1. Prove the exact evaluation}
Put $G=\Gamma(1/4)$. We prove Xu--Zhao's Conjecture 5.5
\cite[p.~24]{XZ}, namely
\begin{align}
 I:=\int_0^\infty\frac{x^9\,dx}{(\cos x+\cosh x)^6}
 ={}&-\frac{63}{5\cdot2^{10}}
 +\frac{1071G^8}{5^2\cdot2^{13}\pi^2}
 -\frac{21G^8}{2^{12}\pi^3}
 +\frac{63G^8}{2^{16}\pi^4}\notag\\
 &-\frac{21G^{16}}{5^3\cdot2^{13}\pi^4}
 +\frac{3G^{16}}{5\cdot2^{13}\pi^5}
 -\frac{161G^{16}}{5\cdot2^{19}\pi^6}
 +\frac{21G^{16}}{2^{21}\pi^8}\notag\\
 &+\frac{G^{24}}{5^2\cdot2^{19}\pi^8}
 -\frac{G^{24}}{3\cdot2^{20}\pi^9}
 +\frac{69G^{24}}{5\cdot2^{25}\pi^{10}}\notag\\
 &-\frac{21G^{24}}{5\cdot2^{24}\pi^{11}}
 +\frac{63G^{24}}{5\cdot2^{27}\pi^{12}}
 -\frac{17G^{32}}{3\cdot5^2\cdot2^{31}\pi^{14}}\notag\\
 &+\frac{13G^{32}}{5\cdot2^{34}\pi^{16}}
 +\frac{3G^{40}}{5^2\cdot2^{40}\pi^{20}}.\label{target}
\end{align}
The source's general structural theorem assumes $p\ge\lfloor m/2\rfloor$
for a numerator $x^{4p+1}$ and denominator power $m$. The present
case $(p,m)=(2,6)$ is outside that assumption. The missing ingredient
is the elementary Lambert-series formula for $S_0$ below.

\subsection*{1. A convergent residue formula}
Let
\[
 F(z)=\frac{z^9}{(\cos z+\cosh z)^6},\qquad
 a=\frac{1+i}{2},\quad b=\frac{1-i}{2}.
\]
The factorization $\cos z+\cosh z=2\cosh(az)\cosh(bz)$ shows
that the only poles in the first quadrant are
\[
 z_n=(1+i)t_n,\qquad t_n=\pi(n+\tfrac12),\quad n\ge0,
\]
each of order six. Integrate on the positively oriented triangle
with vertices $0,R,iR$, where $R=2\pi N$ and $N$ is a positive integer.
On its oblique side $x+y=R$, we have
\[
 |\cosh(az)|=\cosh((x-y)/2)\ge1,\qquad
 |\cosh(bz)|\ge\sinh(R/2).
\]
The modulus of the integral on that side is therefore at most
$\sqrt2R^{10}/(64\sinh^6(R/2))$, which tends to zero.
Since $F(iy)=iF(y)$, the two axial sides together give
$2\int_0^R F(x)\,dx$. The residue theorem consequently yields
\begin{equation}\label{residue-sum}
 I=\pi i\sum_{n\ge0}\Res_{z=z_n}F(z).
\end{equation}
The original integral converges at zero and infinity. The residue
series below is absolutely convergent, because each term is a finite
sum of polynomials in $t_n$ times derivatives of $\sech^6t_n$.

Write $z=z_n+w$. Then
\[
 F(z_n+w)=-\frac{(z_n+w)^9}{64\sinh^6(aw)\cosh^6(t_n+bw)},
 \qquad
 \left(\frac{v}{\sinh v}\right)^6
       =1-v^2+\frac8{15}v^4+O(v^6).
\]
Extracting the coefficient of $w^{-1}$ gives
\begin{equation}\label{residue}
 \Res_{z=z_n}F(z)=\frac i2\sum_{h=0}^2\gamma_h
   \sum_{\ell=0}^{5-2h}\binom9\ell
        \frac{t_n^{9-\ell}}{r!}
          (\sech^6)^{(r)}(t_n),
 \quad r=5-2h-\ell,
\end{equation}
where $(\gamma_0,\gamma_1,\gamma_2)=(1,1,8/15)$.
For example, if $(A_0,A_1,A_2)=(1,-1,8/15)$, the phase of the
term indexed by $h,\ell,r$ is
\[
 -\frac{i}{8}A_h a^{2h}(1+i)^{9-\ell}b^r
          =\frac i2\gamma_h,
 \qquad \ell+r+2h=5.
\]
This checks both the sign and all powers of two in \eqref{residue}.

Define, for even nonnegative $m$ and $y>0$,
\[
 S_m(y)=\sum_{n\ge0}(n+\tfrac12)^m
                         \sech^2((n+\tfrac12)y).
\]
The elementary differential identity
\[
 \sech^6t=\frac1{120}
       \left(\frac{d^4}{dt^4}-20\frac{d^2}{dt^2}+64\right)\sech^2t
\]
and \eqref{residue-sum}--\eqref{residue} now imply the finite formula
\begin{equation}\label{finite}
 I=-\frac\pi{240}\sum_{h=0}^2\gamma_h
       \sum_{\ell=0}^{5-2h}\binom9\ell\frac{\pi^{9-\ell}}{r!}
       \sum_{j\in\{0,2,4\}}c_j
                     S_{4+2h-j}^{(r+j)}(\pi),
 \quad r=5-2h-\ell,
\end{equation}
where $(c_0,c_2,c_4)=(64,-20,1)$.
Only $S_0,S_2,S_4,S_6,S_8$ occur. Differentiation under each sum
is justified by locally uniform exponential convergence for $y>0$.

\subsection*{2. The five sums as Lambert series}
For $|q|<1$ and positive odd $s$, set
\[
 L_s(q)=\sum_{k\ge1}\frac{k^s q^k}{1-q^k},\qquad D=q\frac d{dq}.
\]
In the next identities $q=e^{-y}$. For $m\in\{2,4,6,8\}$,
\begin{equation}\label{Sm}
 S_m(y)=2^{2-m}D\left\{L_{m-1}(q)
       -(2+2^{m-1})L_{m-1}(q^2)+2^mL_{m-1}(q^4)\right\}.
\end{equation}
Indeed, $\sech^2(ky/2)=4q^k/(1+q^k)^2$, and
\[
 \sum_{\substack{k\ge1\\k\text{ odd}}}
       \frac{k^{m-1}q^k}{1+q^k}
 =L_{m-1}(q)-(2+2^{m-1})L_{m-1}(q^2)
                        +2^mL_{m-1}(q^4).
\]
Differentiating proves \eqref{Sm}.

For the remaining case, interchange two absolutely convergent sums:
\begin{align*}
 \frac14 S_0(y)
 &=\sum_{\substack{k\ge1\\k\text{ odd}}}
       \sum_{v\ge1}(-1)^{v-1}v q^{kv}
   =\sum_{v\ge1}\frac{(-1)^{v-1}v q^v}{1-q^{2v}}\\
 &=L_1(q)-5L_1(q^2)+4L_1(q^4).
\end{align*}
Thus
\begin{equation}\label{Szero}
 S_0(y)=4\{L_1(q)-5L_1(q^2)+4L_1(q^4)\}.
\end{equation}
This uses a weight-two Lambert series instead of a negative-index
series, and is valid for every $y>0$.

\subsection*{3. Exact modular values and derivatives}
Use the standard normalization $q=e^{2\pi i\tau}$ for Eisenstein
series. In this normalization,
\[
 L_1=\frac{1-E_2}{24},\qquad
 L_3=\frac{E_4-1}{240},\qquad
 L_5=\frac{1-E_6}{504},\qquad
 L_7=\frac{E_4^2-1}{480}.
\]
Ramanujan's differential identities \cite[Section 7]{R} are
\begin{equation}\label{ramanujan}
 DE_2=\frac{E_2^2-E_4}{12},\qquad
 DE_4=\frac{E_2E_4-E_6}{3},\qquad
 DE_6=\frac{E_2E_6-E_4^2}{2}.
\end{equation}
They give every derivative needed in \eqref{finite} by polynomial
arithmetic over $\Q$.

Put
\[
 u=\theta_3(0\mid i)^4=\frac{G^4}{4\pi^3}.
\]
The last equality is the classical elliptic-integral value also
recorded in \cite[equation (2.3)]{XZ}. The three evaluation points
in \eqref{Sm}--\eqref{Szero} give
\begin{equation}\label{values}
\begin{array}{c|ccc}
\tau&E_2(\tau)&E_4(\tau)&E_6(\tau)\\ \hline
i/2&6/\pi-3u/2&33u^2/4&-189u^3/8\\
i&3/\pi&3u^2/4&0\\
2i&3/(2\pi)+3u/8&33u^2/64&189u^3/512
\end{array}
\end{equation}
For completeness these values follow from the classical theta
identities as follows. At $i$ the fourth powers of $\theta_2$ and
$\theta_4$ are both $u/2$. Theta duplication \cite{DLMF} gives at $2i$
\[
 (\theta_2^4,\theta_3^4,\theta_4^4)
   =u\left(\frac{3-2\sqrt2}{8},
           \frac{3+2\sqrt2}{8},\frac1{\sqrt2}\right).
\]
Substitute into
\[
 E_4=\tfrac12(\theta_2^8+\theta_3^8+\theta_4^8),\qquad
 E_6=\tfrac12(\theta_2^4+\theta_3^4)
                  (\theta_3^4+\theta_4^4)(\theta_4^4-\theta_2^4).
\]
For weight two use
\[
 2E_2(2\tau)-E_2(\tau)=\tfrac12(\theta_3^4+\theta_4^4),\qquad
 E_2(-1/\tau)=\tau^2E_2(\tau)+\frac{6\tau}{\pi i}.
\]
The latter gives $E_2(i)=3/\pi$. These formulas and the modular
transformations $E_w(-1/\tau)=\tau^wE_w(\tau)$ for $w=4,6$
give all of \eqref{values}.

Here are explicit instructions for the finite evaluation, including
the chain factors. Introduce variables $e_2,e_4,e_6$ and the derivation
\[
 \mathcal D=\frac{e_2^2-e_4}{12}\partial_{e_2}
       +\frac{e_2e_4-e_6}{3}\partial_{e_4}
       +\frac{e_2e_6-e_4^2}{2}\partial_{e_6}.
\]
Let $\mathcal L_w=(1-e_2)/24,(e_4-1)/240,(1-e_6)/504,
(e_4^2-1)/480$ for $w=2,4,6,8$, respectively. Let $V_a$ denote
substitution from row $\tau=ai/2$ of \eqref{values}, for $a=1,2,4$.
For $d\ge0$,
\begin{align}
 S_0^{(d)}(\pi)
   &=4(-1)^d\sum_{a=1,2,4} b_a a^d
                         V_a(\mathcal D^d\mathcal L_2),
       &(b_1,b_2,b_4)&=(1,-5,4),\label{evalzero}\\
 S_m^{(d)}(\pi)
   &=2^{2-m}(-1)^d\sum_{a=1,2,4} b_{m,a}a^{d+1}
                         V_a(\mathcal D^{d+1}\mathcal L_m),
       &&m=2,4,6,8,\label{evalpositive}
\end{align}
where $(b_{m,1},b_{m,2},b_{m,4})=(1,-2-2^{m-1},2^m)$.
Here the factor $(-1)^d$ comes from $d/dy=-D$, and each factor
$a$ comes from differentiating $q^a$. Differentiation is performed
\emph{before} substituting the special values.

\subsection*{4. The resulting polynomial}
Substitution of \eqref{evalzero}--\eqref{evalpositive} into
\eqref{finite}, followed by rational polynomial arithmetic, gives
\begin{align*}
 I={}&-\frac{63}{5120}
 +u^2\left(\frac{1071\pi^4}{12800}-\frac{21\pi^3}{256}
                                      +\frac{63\pi^2}{4096}\right)\\
 &+u^4\left(-\frac{21\pi^8}{4000}+\frac{3\pi^7}{160}
                   -\frac{161\pi^6}{10240}+\frac{21\pi^4}{8192}\right)\\
 &+u^6\left(\frac{\pi^{10}}{3200}-\frac{\pi^9}{768}
       +\frac{69\pi^8}{40960}-\frac{21\pi^7}{20480}
       +\frac{63\pi^6}{163840}\right)\\
 &+u^8\left(-\frac{17\pi^{10}}{2457600}
                         +\frac{13\pi^8}{1310720}\right)
 +\frac{3\pi^{10}u^{10}}{26214400}.
\end{align*}
Replacing $u$ by $G^4/(4\pi^3)$ is precisely \eqref{target}.
The accompanying exact checker implements the displayed derivation
and substitutions and verifies zero difference from every one of the
sixteen target terms. Its certificate records the nine grouped
contributions indexed by $(h,j)$ in \eqref{finite}, so the arithmetic
can also be checked group by group. No numerical recognition of
constants or unproved convergence assertion is used.\qed

\section*{Task 2. Remove the low-degree restriction}
Extend this argument to $\int_0^\infty x^{4p+1}(\cos x+\cosh x)^{-m}\,dx$
when $p<\lfloor m/2\rfloor$. Determine which negative-index sums
produced by the residue calculation can be replaced by positive-weight
Lambert series through interchange of divisor sums.

\section*{Task 3. Prove the fifth-moment pattern}
Apply the resulting machinery to Conjecture 5.9 of \cite{XZ}, concerning
all denominator powers with numerator $x^5$. Explain the exact allowed
powers of $\Gamma(1/4)$ and $\pi$, including the predicted missing
coefficient, rather than merely producing individual evaluations.

\section*{Task 4. Evaluate at other modular parameters}
Replace the equal cosine and hyperbolic-cosine scales by a ratio tied
to a general elliptic modulus. Identify analogous exact evaluations
at other imaginary quadratic parameters and determine the necessary
periods and algebraic coefficients.

\section*{Task 5. Find recurrences for exact evaluation}
Develop recurrences in numerator degree and denominator power for the
residue polynomials and their modular expressions. Measure the growth
of intermediate expressions and identify a basis that makes the final
cancellations visible before expansion.

\begin{thebibliography}{3}
\bibitem{XZ} C.~Xu and J.~Zhao,
\emph{General Berndt-Type Integrals and Series Associated with Jacobi
Elliptic Functions},
\href{https://arxiv.org/abs/2401.01385v2}{arXiv:2401.01385v2} (2024).
\bibitem{R} S.~Ramanujan, \emph{On certain arithmetical functions},
Trans. Cambridge Philos. Soc. \textbf{22} (1916), 159--184,
Section 7; \href{https://ramanujan.sirinudi.org/Volumes/published/ram18.html}{online transcription}.
\bibitem{DLMF} NIST Digital Library of Mathematical Functions,
\href{https://dlmf.nist.gov/20.7}{Section 20.7}, theta identities and
transformations, and \href{https://dlmf.nist.gov/23.8}{Section 23.8},
elliptic-function series.
\end{thebibliography}
\noindent\textit{Verification and status.}
Run \texttt{verify\_evaluation.py} with Python, SymPy, and mpmath.
Exact algebra verifies the finite certificate. The separate numerical
checks are diagnostics only. The full current primary source retains
Conjecture 5.5; dated source and claim audits accompany this note.
The authoring agent performed the proof review, without independent
peer review.
\end{document}
